\rtmtight
\begin{tblr}{width=\textwidth,colspec={|Q[c,m,2.1cm]|X[l,m]|},hlines,vlines,hspan=minimal,rowsep=1.5pt,colsep=3pt,row{1}={bg=headgray,font=\bfseries}}
\SetCell{c,m}\hfil\logo\hfil & \textbf{POLITEKNIK NEGERI MALANG}\par\textbf{JURUSAN TEKNOLOGI INFORMASI}\par\textbf{PROGRAM STUDI: D4 TEKNIK INFORMATIKA} \\
\SetCell[c=2]{l} \textbf{RENCANA TUGAS MAHASISWA (RTM) DAN RUBRIK PENILAIAN} & \\
Nama Mata Kuliah & Bahasa Inggris Persiapan Kerja \\
Kode Mata Kuliah & RTI258002 \quad \textbf{sks / Jam perminggu:} 2 SKS/4 jam \quad \textbf{Semester:} 8 \\
Dosen Pengampu & \begin{enumerate}\item Farid Angga Pribadi, S.Kom., M.Kom.\end{enumerate} \\
\end{tblr}
\assessmentblock{Tugas Dokumen Profesional}{Case Method dan portofolio}{SCPMK0303-05901, SCPMK0102-05902}{Mahasiswa menyusun portofolio dokumen profesional Bahasa Inggris secara bertahap, mencakup CV, LinkedIn, cover letter, email profesional, dan dokumentasi teknis TI, yang dievaluasi melalui tugas, kuis, dan UTS.}{Menganalisis contoh dokumen profesional, menyusun dokumen sesuai konteks TI, mendapatkan umpan balik peer dan instruktur, serta memperbaiki dokumen berulang.}{Portofolio dokumen profesional (CV, cover letter, email, technical writing), jawaban kuis tertulis, dan presentasi UTS portofolio.}{Kualitas CV, cover letter, dan email profesional & 18 & CV menampilkan pengalaman TI secara strategis, cover letter disesuaikan dengan posisi, dan email menggunakan konvensi profesional yang tepat; semuanya bebas kesalahan gramatikal. & Dokumen lengkap dan informatif, tetapi ada beberapa kesalahan gramatikal atau penyesuaian konteks yang kurang optimal. & Dokumen tidak lengkap, banyak kesalahan gramatikal, atau tidak mengikuti konvensi profesional standar. \\ \\
Kualitas technical writing dan dokumentasi TI & 12 & Dokumentasi teknis tersusun logis, menggunakan terminologi TI yang tepat, ditujukan pada audiens yang tepat, dan disertai contoh/diagram yang relevan. & Technical writing cukup jelas dan menggunakan terminologi yang benar, tetapi struktur atau kesesuaian audiens masih bisa ditingkatkan. & Technical writing tidak terstruktur, terminologi tidak tepat, atau tidak sesuai dengan konteks TI. \\ \\
Penguasaan konsep dokumen profesional (Kuis \& UTS) & 20 & Kuis dan UTS menunjukkan pemahaman mendalam tentang konvensi dokumen profesional Bahasa Inggris dalam konteks TI, dengan contoh dan aplikasi yang tepat. & Sebagian besar konsep dipahami dengan baik, tetapi beberapa aspek konvensi atau penerapannya dalam konteks TI masih kurang. & Konsep dokumen profesional tidak dipahami, banyak jawaban kuis/UTS keliru, atau tidak dapat menerapkan dalam konteks TI. \\}
\assessmentblock{PBL Mock Interview dan Presentasi Teknis}{Project Based Learning dan simulasi}{SCPMK0102-05902, SCPMK0303-05901}{Mahasiswa menjalani simulasi wawancara kerja dan presentasi teknis dalam Bahasa Inggris, menerapkan teknik STAR method, behavioral interview, dan komunikasi multikultural yang dipelajari.}{Berlatih mock interview berpasangan, menyiapkan jawaban STAR, melakukan presentasi teknis proyek TI, mendapatkan feedback instruktur dan teman sebaya.}{Rekaman video mock interview, slide presentasi teknis proyek TI, jawaban kuis tertulis interview skills, dan refleksi mandiri.}{Penguasaan teknik wawancara kerja (STAR \& behavioral) & 10 & Mahasiswa menerapkan STAR method secara konsisten, memberikan jawaban behavioral yang spesifik dengan contoh dari pengalaman TI, dan merespons pertanyaan teknis dengan tepat. & STAR method diterapkan pada sebagian besar jawaban, tetapi beberapa jawaban kurang spesifik atau contoh tidak relevan dengan konteks TI. & Tidak menerapkan STAR method, jawaban sangat umum, atau tidak dapat merespons pertanyaan interview dengan efektif. \\ \\
Kemampuan presentasi teknis dalam Bahasa Inggris & 9 & Presentasi teknis disampaikan dengan jelas, terstruktur, dan percaya diri; terminologi TI digunakan dengan benar; Q\&A dijawab dengan baik. & Presentasi cukup jelas dan informatif, tetapi ada beberapa ketidakakuratan terminologi atau kurang percaya diri saat Q\&A. & Presentasi tidak terstruktur, terminologi sering salah, atau tidak dapat menjawab pertanyaan teknis. \\ \\
Komunikasi lintas budaya dan profesionalisme & 8 & Mahasiswa menunjukkan kesadaran budaya, menggunakan register bahasa yang tepat untuk konteks profesional global, dan berinteraksi dengan etika kerja yang baik. & Komunikasi cukup profesional tetapi kadang menggunakan register yang tidak tepat atau kurang menunjukkan kesadaran budaya. & Komunikasi tidak profesional, register bahasa tidak sesuai, atau tidak menunjukkan pemahaman etika kerja global. \\}
\assessmentblock{UAS Final Mock Interview Panel dan Presentasi Portofolio}{UAS praktik}{SCPMK0102-05902, SCPMK0303-05901}{Mahasiswa mengikuti sesi final berupa mock interview panel dan presentasi portofolio profesional di depan tim instruktur sebagai simulasi proses rekrutmen TI yang sesungguhnya.}{Menyiapkan portofolio final, menjalani mock interview panel (30 menit), mempresentasikan portofolio karier, dan merespons pertanyaan panel.}{Portofolio profesional final (digital), rekaman/live mock interview panel, dan slide presentasi portofolio.}{Performa mock interview panel & 10 & Mahasiswa menjawab semua pertanyaan panel (HR, manajer teknis, dan penguji) dengan lancar, tepat, dan profesional menggunakan Bahasa Inggris baku. & Sebagian besar pertanyaan dijawab dengan baik, tetapi ada beberapa pertanyaan teknis atau behavioral yang kurang memuaskan. & Tidak dapat menjawab banyak pertanyaan panel, bahasa tidak profesional, atau tidak siap menghadapi format panel interview. \\ \\
Kualitas presentasi portofolio profesional & 8 & Portofolio lengkap, profesional, dan dipresentasikan secara meyakinkan; mahasiswa dapat mengartikulasikan nilai jual dirinya dalam Bahasa Inggris. & Portofolio lengkap dan presentasi cukup baik, tetapi artikulasi nilai jual atau kelancaran penyampaian masih dapat ditingkatkan. & Portofolio tidak lengkap atau presentasi tidak meyakinkan; mahasiswa tidak dapat menjelaskan kompetensinya secara efektif. \\ \\
Kelancaran dan ketepatan Bahasa Inggris & 5 & Bahasa Inggris digunakan dengan lancar, akurat secara gramatikal, dan sesuai konteks profesional sepanjang sesi UAS. & Bahasa Inggris umumnya dapat dipahami dengan baik, tetapi ada beberapa kesalahan gramatikal atau penggunaan yang kurang tepat. & Banyak kesalahan gramatikal yang mengganggu komunikasi, atau kosakata tidak memadai untuk konteks profesional. \\}

\begin{tblr}{width=\textwidth,colspec={|Q[l,m,3.2cm]|X[l,m]|},hlines,vlines,hspan=minimal,row{1-3}={bg=headgray},rowsep=1.5pt,colsep=3pt}
Jadwal Pelaksanaan: & Minggu 1--17 sesuai RPS. Setiap evaluasi memiliki RTM dan rubrik multi-kriteria. \\
Lain-Lain Yang Diperlukan: & LMS, portofolio digital, aplikasi presentasi, rekam video, dan perangkat presentasi. \\
Pustaka: & Mengacu pustaka utama dan pendukung pada bagian RPS. \\
\end{tblr}
